
%background
\documentclass[main.tex]{subfiles}

\begin{document}

\section{Background on Jouanolou models}\label{s:background}
\subsection{Definitions and preliminary results}

Refer to \cite[\S 1.2]{FHK}, \cite[\S 1]{GWWchiral}, and specifically \cite[\S 1]{GWvir2d} for more details.

\paragraph
Let $\cJ_d^{poly}$ denote the Jouanlou model for the structure sheaf on $\AA^d - \{0\}$.
It is a dg algebra defined over $\C[z] = \C[z^\1, \ldots, z^\dd]$.
As a graded ring it is the following quotient of a graded polynomial ring
\begin{equation}\label{}
  \frac{\C[z^\s,x^\s,\d x^\s]_{\1 \leq \s \leq \dd}}{\<\sum_{\s=\1}^d z^\s x^\s -1, \quad \sum_{\s=\1}^d z^\s \d x^\s \> }
\end{equation}
where $z^\s, x^\s$ carry degree zero, and $\d x^\s$ carries degree $+1$ for each $\s$.
Note that $\cJ_d^{poly}$ is concentrated in degrees $0,\ldots, d-1$.
The differential is defined by
\[
  \dbar z^\s = 0, \quad \dbar x^\s = \d x^\s .
\]
Let $\cJ_d$ be its completion over the $z$-variables.
The latter is a dg algebra defined over $\C\llbracket z \rrbracket = \C\llbracket z^\1, \ldots, z^\dd \rrbracket$.

\paragraph
The Jouanolou model $\cJ_d$ (and $\cJ_d^{poly}$) is a $GL_d$-equivariant dg algebra meaning the product and
differential are $GL_d$
equivariant.
Let $\lie{q}$ denote the $d$-dimensional defining representation of $GL_d$.

\begin{proposition}[\cite{FHK}]
\label{prop:Jouanolou}
The cohomology of $\cJ^{poly}_d$ is concentrated in degrees zero and $d-1$.
There are natural $GL_d$-equivariant isomorphisms
\begin{equation}\label{}
  H^0(\cJ_d^{poly}) \cong \op{S}^\bu (\lie{q}^*) = \C[z]
\end{equation}
and
\begin{equation}\label{}
  H^{d-1} (\cJ_d^{poly}) \cong \op{S}^\bu (\lie{q}) \otimes \op{det}(\lie{q}) .
\end{equation}
Similar results hold for $\cJ_d$ where the first isomorphism is replaced by
\[
  H^0(\cJ_d) \cong \Hat{\op{S}}(\lie{q}^*) = \C\llbracket z \rrbracket .
\]
\end{proposition}

Let
\begin{equation}\label{}
  P \define \sum_{\s=1}^d (-1)^{\s-1} x^{\s} \d x^\1 \cdots \Hat{\d x^\s} \cdots \d x^\dd
\end{equation}
where the hat means that we exclude that wedge factor.
This is an element of $\cJ_d$ of degree $(d-1)$ and transforms in the determinant representation of $GL_d$.
For $\bk = (k_\1, \ldots, k_\dd)$ a $d$-tuple of non-negative integers,
\[
\del^{(\bk)} P \define \frac{(\sum_\ii k_\ii +1)!}{\prod_\ii k_\ii} (x^\1)^{k_\1} \cdots (x^\dd)^{k_\dd}  P .
\]
The collection $\{[\del^{(\bk)} P]_{\dbar}\}$ where $\bk$ ranges over all $d$-tuples of non-negative integers is a basis for
$H^{d-1}(\cJ_d)$.

\paragraph
The Jouanolou model for the tangent sheaf $\sT$ on $\AA^{d} - \{0\}$ is the polynomial dg Witt algebra denoted
$\lie{witt}_d^{poly}$.
It is naturally a dg Lie algebra whose differential we will also denote $\dbar$.
It is freely generated over $\cJ_d^{poly}$ by generators we denote $\del_\1, \ldots, \del_\dd$.
Equivalently, as a
cochain complex
\begin{equation}\label{}
  \lie{witt}^{poly}_d = \cJ_d^{poly} \otimes \lie{q}
\end{equation}
where $\lie{q} = \op{span}\{\del_\1, \ldots, \del_{\dd}\}$.

There is a natural $\cJ_d$-linear quasi-isomorphism of dg Lie algebras
\begin{equation}\label{}
  \lie{witt}_d^{poly} \xto{\simeq} \op{Der} \cJ^{poly}_d
\end{equation}
which is defined by
\[
  \del_{\ii} z^\jj = \delta_{\ii}^{\jj}, \quad \del_{\ii} x^{\jj} = -x^\ii x^\jj, \quad \del_{\ii} \d x^{\jj} = -
  x^{\jj} \d x^{\ii} - x^{\ii} \d x^{\jj} .
\]
As is usual, if $X \in\lie{witt}_d^{poly}$ then we denote $X f$ its action on $f \in \cJ_d$.
Similar remarks and formulas are true for the formal completion over the $z$-variables denoted $\lie{witt}_d$.

\begin{proposition}[\cite{GWvir2d}]
\label{prop:JouanolouWitt}
The cohomology of $\lie{witt}^{poly}_d$ is concentrated in degrees zero and $d-1$.
There is a natural $GL_d$-equivariant isomorphism of Lie algebras
\begin{equation}\label{}
  H^0(\lie{witt}_d^{poly}) \cong \op{S}^\bu (\lie{q}^*) \otimes \lie{q} = \lie{w}_d^{poly}
\end{equation}
and
\begin{equation}\label{}
  H^{d-1} (\lie{witt}_d^{poly}) \cong \op{S}^\bu (\lie{q}) \otimes \op{det}(\lie{q}) \otimes \lie{q}.
\end{equation}
Similar results hold for $\cJ_d$ where the first isomorphism is replaced by
\[
  H^0(\lie{witt}_d) \cong \lie{w}_d .
\]
\end{proposition}


\paragraph{Tensor modules}\label{TensorMod}
Let $V$ be a finite-dimensional $\lie{gl}_d$-module.
Define the cochain complex
\begin{equation}\label{}
  \cV \define \cJ_d \otimes V .
\end{equation}
We extend the representation $\rho_V \colon \lie{gl}_d \to \op{End}(V)$ in a $\cJ_d$-linear way which we denote by the
same symbol
\begin{equation}\label{}
  \rho_V \colon \cJ_d \otimes \lie{gl}_d \to \op{End}(\cV) .
\end{equation}
For $X \in \lie{witt}_d$, define $\mathsf{L}_{V}(X) \in \op{End}(\cV)$ by the formula
\begin{equation}\label{}
  \mathsf{L}_V(X) (f \otimes v) \define (Xf) \otimes v + \rho_V(JX) (f \otimes v)
\end{equation}
where $JX \in \cJ_d \otimes \lie{gl}_d$ is the Jacobian associated to $X$.
If $X = g^\ii \del_{\ii}$ it is defined by
\begin{equation}\label{eq:jacobian}
  (J X )_{\jj}^\ii = \frac{\del g^\ii}{\del z^\jj} \in \cJ_d .
\end{equation}
It is immediate to verify that $X \mapsto \sfL_X$ endows $\cV$ with the structure of a dg $\lie{witt}_d$-module.

We note that $H^0 (\lie{\cV})$ is naturally a \textit{tensor} $H^0(\lie{witt}_d) = \lie{w}_d$-module which is obtained
by restricting $V$ along
\begin{equation}\label{}
  \{X \in \lie{w}_d\; | \; X(0) = 0\} = \lie{w}_d^{\geq 0} \twoheadrightarrow \lie{gl}_d
\end{equation}
and coinducing along $\lie{w}_d^{\geq 0} \hookrightarrow \lie{w}_d$.
For this reason, we refer to such $\cV$ as a tensor module for $\lie{witt}_d$.
These modules appear in our formulation of the universal Grothendieck--Riemann--Roch theorem in section \ref{s:GRR}.

\paragraph{Matrix extension}
For a positive integer $r$ we will consider the dg algebra of $\cJ_d$-valued $r \times r$ matrices $\lie{gl}
_r(\cJ_d)$.
Its differential is induced from $\dbar$ and the product is defined using the commutative multiplication on $\cJ_d$
together with ordinary matrix multiplication.
Via commutator, we may treat it as a dg Lie algebra.
Invariant polynomials on $\lie{gl}_r$ of degree $d+1$ define central extensions of $\lie{gl}_r(\cJ_d)$ called
higher-dimensional Kac--Moody algebras \cite{FHK}.
We are ultimately interested in central extensions of this dg Lie algebra extended by the dg Witt algebra.

Entry-wise, $\lie{witt}_d$ acts on $\lie{gl}_r (\cJ_d)$ by derivations.
We therefore can form the semi-direct product dg Lie algebra
\begin{equation}\label{}
  \lie{witt}_d \ltimes \lie{gl}_r(\cJ_d) .
\end{equation}
In sections \ref{s:chern} and \ref{s:gf} we provide a classification of central extensions of this dg Lie algebra.
This dg Lie algebra is also the main ingredient in our formulation of the universal Grothendieck--Riemann--Roch theorem
in section \ref{s:GRR}.

\subsection{Flatness of the Jouanolou model}
We will use the following algebraic result in the main part of the paper.

\begin{proposition}\label{prop:flatness}
The Jouanolou model $\cJ_d$ is flat over $\C\llbracket z\rrbracket=\C\llbracket z^{\1},\dots,z^{\dd}\rrbracket$. An analogous statement is valid in the polynomial setting.
\end{proposition}

\begin{proof}
  The proof of \cite{FGY} can be modified to apply to the Jouanolou model in this setting. For the reader’s convenience, we present a more direct argument.

  As above, we denote $\C\llbracket z \rrbracket = \C\llbracket z^\1, \ldots, z^\dd \rrbracket$ and $\C[x] = \C[x^\1,
\ldots,x^\dd]$.
Consider the degree-zero component
\[
\cJ_d^0=\frac{\C\llbracket z\rrbracket[x]}{\<\sum_{\s=1}^d z^{\s}x^{\s}
-1\>}
\]
which is naturally an algebra over $\C\llbracket z \rrbracket$. We will show that $\cJ_d^0$ is flat as a $\C\llbracket z\rrbracket$-algebra, and that $\cJ_d$ is flat over $\cJ_d^0$.

Let $\mathbf{p}$ be a prime ideal in $\cJ_d^0$, and denote $\mathbf{q}=\mathbf{p}\cap \C\llbracket z \rrbracket$. Since $\sum_{\s=1}^d z^{\s}x^{\s}=1$ in $\cJ_d^0$, there exists $\ii$ such that $z^{\ii}\notin \mathbf{q}$. Consequently, $z^{\ii}$ becomes invertible in the localizations $\C\llbracket z \rrbracket_{\mathbf{q}}$ and $(\cJ_d^0)_{\mathbf{p}}$. Observe that
\[
\cJ_d^0\big[\tfrac{1}{z^{\ii}}\big]\simeq
\C\llbracket z \rrbracket\big[\tfrac{1}{z^{\ii}}\big][x^{\1},\dots,\widehat{x^{\ii}},\dots,x^{\dd}],
\]
where $\widehat{x^{\ii}}$ indicates that the variable $x^{\ii}$ is omitted. Localizing further at $\mathbf{p}$, we obtain
\[
(\cJ_d^0)_{\mathbf{p}}\simeq
\big(\C\llbracket z \rrbracket_{\mathbf{q}}[x^{\1},\dots,\widehat{x^{\ii}},\dots,x^{\dd}]\big)_{\mathbf{p}'},
\]
where $\mathbf{p}'=(\C\llbracket z \rrbracket-\mathbf{q})^{-1}\mathbf{p}$. Hence $(\cJ_d^0)_{\mathbf{p}}$ is flat over $\C\llbracket z \rrbracket_{\mathbf{q}}$ for every prime $\mathbf{p}\subset \cJ_d^0$. It follows that $\cJ_d^0$ is flat over $\C\llbracket z \rrbracket$.

Observe that we have the following short exact sequence of $\cJ^0_d$-modules:
\[
0\rightarrow \cJ^0_d\xrightarrow{\iota} (\cJ^0_d)^{\oplus d}\rightarrow \cJ^1_d\rightarrow 0,
\]
where $\iota(f)=f\bigl(\sum_{\s=1}^d z^{\s}\bar{\partial}x^{\s}\bigr)$. Consider the map $r\colon (\cJ^0_d)^{\oplus d}\rightarrow \cJ^0_d$ defined by
\[
r(f_1,\dots,f_d)=\sum_{\s=1}^d f_{\s}x^{\s}.
\]
Then $r\circ \iota=\mathrm{id}$, which shows that $\cJ^1_d$ is a direct summand of $(\cJ^0_d)^{\oplus d}$. Any
direct summand of a free $\cJ^0_d$-module is projective, hence in particular flat. Consequently, $\cJ_d=\wedge^\bu_{\cJ^0}\cJ^1_d$ is flat over $\cJ^0_d$, and therefore also flat over $\C\llbracket z \rrbracket$.
\end{proof}

\end{document}
